\documentclass[11pt]{article}
% ===========================================================================
%  Shared preamble for the AP Calculus AB exam-preparation worksheets.
%  Input right after \documentclass. Files beginning with "_" are not compiled
%  on their own. Navy and gold, to match the course's lesson pages.
% ===========================================================================
\usepackage[letterpaper,top=2.3cm,bottom=2.2cm,left=1.7cm,right=1.7cm,headheight=28pt,footskip=30pt]{geometry}
\usepackage{amsmath,amssymb}
\usepackage[T1]{fontenc}
\usepackage[utf8]{inputenc}
\usepackage{lmodern}
\usepackage{xcolor}
\usepackage[most]{tcolorbox}
\usepackage{enumitem}
\usepackage{fancyhdr}
\usepackage{array}
\usepackage{tikz}
\usepackage{pgfplots}
\pgfplotsset{compat=1.18}

\definecolor{navy}{HTML}{161E54}
\definecolor{ink}{HTML}{27357E}
\definecolor{gold}{HTML}{B98A00}
\definecolor{goldbg}{HTML}{FFFDF3}
\definecolor{inkbg}{HTML}{F6F7FD}
\definecolor{rule}{HTML}{C9CFF0}
\definecolor{okgreen}{HTML}{1B7A44}
\definecolor{okbg}{HTML}{F1F8F3}

\pagestyle{fancy}
\fancyhf{}
\fancyhead[L]{\footnotesize Name: \underline{\hspace{4cm}}}
\fancyhead[C]{\textbf{\color{navy}Integration Academy}\\[-2pt]{\scriptsize Calculus AB \textperiodcentered{} AP\textsuperscript{\textregistered} Exam Preparation}}
\fancyfoot[L]{\tiny\color{black!55}\parbox[t]{0.78\textwidth}{AP\textsuperscript{\textregistered} is a registered trademark of the College Board, which was not involved in the production of, and does not endorse, this material. \textcopyright{} Integration Academy.}}
\fancyfoot[R]{\scriptsize\thepage}
\renewcommand{\headrulewidth}{1.2pt}
\renewcommand{\headrule}{\hbox to\headwidth{\color{gold}\leaders\hrule height \headrulewidth\hfill}}
\renewcommand{\footrulewidth}{0pt}
\setlength{\parindent}{0pt}
\setlength{\parskip}{4pt}

% Title block: \wstitle{code}{title}{unit line}
\newcommand{\wstitle}[3]{\begin{center}{\small\bfseries\color{gold}WORKSHEET #1}\\[3pt]{\LARGE\bfseries\color{navy}#2}\\[3pt]{\small\color{black!60}#3}\end{center}\vspace{2pt}}
% Section heading
\newcommand{\sech}[1]{\par\vspace{10pt}{\large\bfseries\color{navy}#1}\par\nobreak\vspace{2pt}{\color{rule}\hrule height0.8pt}\vspace{7pt}}
% A level tag printed beside a group of questions
\newcommand{\level}[1]{\par\vspace{6pt}{\footnotesize\bfseries\color{gold}\MakeUppercase{#1}}\par\nobreak\vspace{2pt}}

\newtcolorbox{ideabox}{colframe=ink,colback=inkbg,coltitle=white,colbacktitle=ink,fonttitle=\bfseries\small,title={Key ideas},arc=2pt,boxrule=0.8pt,left=8pt,right=8pt}
\newtcolorbox{example}[1]{colframe=rule,colback=white,coltitle=navy,colbacktitle=inkbg,fonttitle=\bfseries\small,title={#1},arc=2pt,boxrule=0.8pt,borderline west={3pt}{0pt}{ink},left=10pt,right=8pt}
\newtcolorbox{frbox}[1]{colframe=gold,colback=goldbg,coltitle=white,colbacktitle=gold,fonttitle=\bfseries\small,title={#1},arc=2pt,boxrule=0.8pt,breakable,left=8pt,right=8pt}
\newtcolorbox{solbox}{colframe=okgreen,colback=okbg,coltitle=white,colbacktitle=okgreen,fonttitle=\bfseries\small,title={Worked solutions},arc=2pt,boxrule=0.8pt,breakable,left=8pt,right=8pt}
\newtcolorbox{scorebox}[1]{colframe=navy,colback=white,coltitle=white,colbacktitle=navy,fonttitle=\bfseries\small,title={#1},arc=2pt,boxrule=0.8pt,breakable,left=8pt,right=8pt}

\newcommand{\soln}{\par\smallskip\textbf{\color{ink}Solution.}\ }
\newcommand{\result}[1]{\par\smallskip\textbf{#1}}
% blank working space of a given height
\newcommand{\work}[1]{\par\vspace{#1}}
% a point value in the margin of a scoring guide
\newcommand{\pts}[1]{\hfill{\footnotesize\bfseries\color{gold}[#1]}}

% Coordinate grid for graph questions: \begin{apgraph}{xmin}{xmax}{ymin}{ymax} ... \end{apgraph}
\newenvironment{apgraph}[4]{\begin{center}\begin{tikzpicture}\begin{axis}[
  axis lines=middle, xmin=#1, xmax=#2, ymin=#3, ymax=#4,
  grid=both, grid style={black!12}, xtick distance=1, ytick distance=1,
  ticklabel style={font=\scriptsize}, xlabel={$x$}, ylabel={$y$},
  width=10cm, height=7.2cm, samples=80, clip=false,
  every axis plot/.append style={line width=1.1pt, color=ink}]}{\end{axis}\end{tikzpicture}\end{center}}
\newcommand{\closeddot}[2]{\addplot[only marks, mark=*, mark size=2.2pt, color=ink] coordinates {(#1,#2)};}
\newcommand{\opendot}[2]{\addplot[only marks, mark=*, mark size=2.2pt, color=ink, mark options={fill=white}] coordinates {(#1,#2)};}

% Multiple-choice options: four across, or stacked when the options are long.
\newcommand{\opts}[4]{\par\vspace{3pt}\begin{tabular}{@{}p{0.235\linewidth}p{0.235\linewidth}p{0.235\linewidth}p{0.235\linewidth}@{}}(A)\ #1 & (B)\ #2 & (C)\ #3 & (D)\ #4\end{tabular}\par\vspace{9pt}}
\newcommand{\optsv}[4]{\par\vspace{3pt}\begin{tabular}{@{}l@{\ \ }p{0.9\linewidth}@{}}(A) & #1\\(B) & #2\\(C) & #3\\(D) & #4\end{tabular}\par\vspace{9pt}}

\fancyhead[R]{\footnotesize\color{black!60}Worksheet 4.1}
\begin{document}

\wstitle{4.1}{What a Derivative Means, in Words and Units}{Unit 4: Derivatives in Context}

\begin{ideabox}
\begin{itemize}[leftmargin=*,itemsep=2pt,topsep=2pt]
  \item $f'(a)$ is the \textbf{rate} at which $f$ is changing at the instant $x=a$. Its units are the units of $f$ per unit of $x$.
  \item A complete interpretation names the instant, the quantity, whether it is increasing or decreasing, the rate (as a positive number) and the units.
  \item If $f$ is itself a rate, $f'$ is a rate of change of a rate: for example litres per minute \emph{per minute}.
  \item A value and a rate give a nearby estimate: $f(a+h)\approx f(a)+f'(a)\,h$.
\end{itemize}
\end{ideabox}

\sech{Worked examples}
\begin{example}{Example 1 \textperiodcentered{} Interpreting}
The cost of producing $q$ kilograms of flour is $C(q)$ dollars. Interpret $C'(40)=3.5$.\soln The units are dollars per kilogram, at the production level $q=40$.\result{When 40 kg are being produced, the cost is increasing at 3.5 dollars per additional kilogram.}
\end{example}

\begin{example}{Example 2 \textperiodcentered{} Estimating}
A plant has height $h(t)$ cm after $t$ weeks, with $h(6)=30$ and $h'(6)=2.5$. Estimate its height at 8 weeks.\soln $h(8)\approx30+2.5(2)$.\result{About 35 cm.}
\end{example}

\begin{example}{Example 3 \textperiodcentered{} A rate of a rate}
A car's velocity is $v(t)$ metres per second at time $t$ seconds. State the units of $v'(3)$ and interpret $v'(3)=-1.2$.\soln Metres per second per second.\result{At $t=3$ seconds the velocity is decreasing at 1.2 metres per second each second.}
\end{example}

\newpage
\sech{Practice}

\level{Core}
\begin{enumerate}[leftmargin=*,itemsep=2pt]
  \item The volume of water in a tank is $V$ litres at time $t$ hours. State the units of $\dfrac{dV}{dt}$.\work{2.3cm}
  \item The temperature of a drink is $T(t)$ degrees Celsius after $t$ minutes. Interpret $T'(15)=-0.8$.\work{2.3cm}
  \item Given $A(20)=150$ and $A'(20)=4$, estimate $A(22)$.\work{2.3cm}
  \item What does $f'(5)<0$ tell you about $f$ at $x=5$?\work{2.3cm}
\end{enumerate}

\level{Exam level}
\begin{enumerate}[leftmargin=*,itemsep=2pt,start=5]
  \item The revenue from selling $x$ items is $R(x)$ dollars. Interpret $R'(100)=1.5$.\work{2.8cm}
  \item A car has $F(t)$ litres of fuel in its tank after $t$ hours of driving. Interpret $F'(2)=-7$.\work{2.8cm}
  \item Gas flows into a cylinder at the rate $G(t)$ litres per minute. State the units of $G'(4)$ and interpret $G'(4)=0.3$.\work{2.8cm}
  \item A patient's weight is $W(t)$ kg after $t$ weeks, with $W(10)=65$, $W'(10)=-0.5$ and $W''(10)=0.02$. Estimate $W(12)$, and say whether the weight loss is speeding up or slowing down at $t=10$.\work{2.8cm}
\end{enumerate}

\begin{frbox}{Question 9 \textperiodcentered{} Free-response practice}
The depth of snow on a field is $S(t)$ centimetres, where $t$ is the number of hours after midnight and $S$ is differentiable with a continuous derivative. It is known that $S(3)=12$, $S'(3)=1.8$ and $S'(9)=-0.6$.
\begin{enumerate}[label=(\alph*),leftmargin=*,itemsep=2pt]
  \item Using correct units, interpret $S'(3)=1.8$.\work{2.6cm}
  \item Use a tangent line to estimate the depth of snow at $t=4$.\work{2.6cm}
  \item Explain why there must be a time between $t=3$ and $t=9$ at which the depth of snow is neither increasing nor decreasing.\work{2.6cm}
\end{enumerate}
\end{frbox}

\newpage
\begin{solbox}
\begin{enumerate}[leftmargin=*,itemsep=4pt]
  \item Litres per hour.
  \item At $t=15$ minutes the temperature of the drink is decreasing at 0.8 degrees Celsius per minute.
  \item $150+4(2)=158$.
  \item $f$ is decreasing at $x=5$.
  \item When 100 items have been sold, revenue is increasing at 1.5 dollars per additional item.
  \item After 2 hours of driving, the amount of fuel in the tank is decreasing at 7 litres per hour.
  \item Litres per minute per minute. At $t=4$ minutes the rate at which gas flows in is increasing by 0.3 litres per minute each minute.
  \item $W(12)\approx65-0.5(2)=64$ kg. The rate $W'$ is negative and increasing (becoming less negative), so the weight loss is slowing down.
  \item \begin{enumerate}[label=(\alph*),leftmargin=*,itemsep=3pt]
    \item At 3 a.m. the depth of snow is increasing at 1.8 centimetres per hour.
    \item $S(4)\approx12+1.8(1)=13.8$ cm.
    \item $S'$ is continuous with $S'(3)>0$ and $S'(9)<0$, so by the Intermediate Value Theorem $S'(c)=0$ for some $c$ between 3 and 9.
  \end{enumerate}
\end{enumerate}
\end{solbox}

\end{document}
